% 7. Conclusion
\section{Conclusion}

In the adult \emph{Drosophila} connectome, neuron-level structural control
impact --- the sensitivity of global directed efficiency to single-neuron
removal --- is not explained by neurotransmitter identity: the
pre-registered GABAergic-excess hypothesis was rejected under
degree-controlled comparison and an ensemble of 100 exact
degree-preserving null networks. Instead, structural chokepoints
concentrate reproducibly in visual-centrifugal architecture, with a
minority of individual neurons --- most notably low-degree centrifugal
bridges --- exceeding degree-matched peers by orders of magnitude. These
results map where the adult fly brain's wiring is most sensitive to
single-neuron perturbation, and provide a ranked, anatomically annotated
target list for functional testing. The analysis is structural
throughout; it identifies network-control-sensitive neurons, not
functionally indispensable ones.
